\begin{pastexam}[label={px:nsysu-112-dmalg-05}]{中山資工所資安組 112 學年度 第 5 題}
（計算）

{[10\%]} Find the generating function for the number of partitions of the nonnegative integer $n$ into
summands, where each summand must appear an odd number of times.
\end{pastexam}

\begin{pastsoln}
答案：
\[
P(x) = \prod_{i \ge 1} \bigl(1 + x^{i} + x^{3i} + x^{5i} + \cdots\bigr)
     = \prod_{i \ge 1} \left(1 + \frac{x^{i}}{1 - x^{2i}}\right).
\]

把 $n$ 拆成正整數的和、不計順序，叫做 $n$ 的一個分割。 一個分割由「每個正整數 $i$ 用了幾次」決定，記
$i$ 用了 $m_i$ 次，則 $\sum_i i \cdot m_i = n$。 題目要求「出現的 summand 都出現奇數次」，也就是每個 $m_i$
是 $0$（沒有用到）或奇數 $1, 3, 5, \ldots$。

\textbf{一個因子管一個 $i$}。 數字 $i$ 用 $m_i$ 次貢獻 $i m_i$ 到總和，用 $x^{i m_i}$ 記錄；$m_i$ 的可能值是
$0, 1, 3, 5, \ldots$，所以 $i$ 的因子是
\[
1 + x^{i} + x^{3i} + x^{5i} + \cdots = 1 + x^{i}\bigl(1 + x^{2i} + x^{4i} + \cdots\bigr) = 1 + \frac{x^i}{1 - x^{2i}}.
\]
把所有 $i \ge 1$ 的因子乘起來，展開時從每個因子各挑一項，挑法對應一組 $(m_1, m_2, \ldots)$，乘積是
$x^{\sum i m_i}$；所以 $x^n$ 的係數就是滿足條件的分割數。 （對固定的 $n$，只有 $i \le n$ 的因子會影響 $x^n$
的係數，乘積當成形式冪級數有意義。）

\textbf{驗算前幾項}。 展開得 $P(x) = 1 + x + x^2 + 3x^3 + 2x^4 + 5x^5 + \cdots$。 例如 $n = 3$：$3$、$2 + 1$、
$1 + 1 + 1$ 三種都符合（$1$ 出現 $3$ 次是奇數）；$n = 4$：$4$、$3 + 1$ 符合，$2 + 2$（$2$ 出現兩次）、$2 + 1 + 1$
（$1$ 出現兩次）、$1 + 1 + 1 + 1$（四次）不符合，共 $2$ 種，和係數一致。
\end{pastsoln}
